% Generated draft: requires mathematical and source review.
\documentclass[11pt]{article}
\usepackage[margin=25mm]{geometry}
\usepackage{fontspec}
\setmainfont{Latin Modern Roman}
\usepackage{amsmath,amssymb,mathtools,unicode-math}
\setmathfont{Latin Modern Math}
\usepackage{enumitem,needspace,xurl,hyperref}
\setlength{\parindent}{0pt}
\setlength{\parskip}{5pt}
\setlength{\emergencystretch}{3em}
\newcommand{\N}{\mathbb{N}}
\newcommand{\Z}{\mathbb{Z}}
\newcommand{\Q}{\mathbb{Q}}
\newcommand{\R}{\mathbb{R}}
\newcommand{\C}{\mathbb{C}}
\newcommand{\F}{\mathbb{F}}
\newcommand{\A}{\mathbb{A}}
\newcommand{\PP}{\mathbb{P}}
\newcommand{\E}{\mathbb{E}}
\newcommand{\eps}{\varepsilon}
\newcommand{\dd}{\,\mathrm{d}}
\begin{document}
\section*{UnsolvedMath ID: 30002829}
\textbf{Source problem code:} catalogue entry
\subsection{Definitions and mathematical statement}
The supplied record does not specify a ground field. Let \(k\) denote the intended ground field; this is a placeholder for missing source data, not a universal quantifier over fields. For \(m\in\{4,6\}\), let \(E_m\) and \(g_m\) denote the elliptic curve and distinguished automorphism intended by the source. The record calls \(E_4\) anharmonic and \(E_6\) equianharmonic, and calls \(g_4,g_6\) their standard automorphisms, without supplying equations or formulas. Assume these data and the quotients below are defined over \(k\).

For a positive integer \(n\), use the diagonal action
\[
 g_m(P_1,\ldots,P_n)=(g_m(P_1),\ldots,g_m(P_n)),
 \qquad Q_{n,m}=E_m^n/\langle g_m\rangle.
\]
The source denotes by \(X_{n,m}\) the minimal resolution of \(Q_{n,m}\). In dimensions four and five, the supplied record does not explain the intended meaning or choice of this minimal resolution. Whenever the intended resolution exists over \(k\), it is birational to \(Q_{n,m}\); the rationality questions therefore depend only on this birational class.

An integral \(d\)-dimensional variety over \(k\) is \(k\)-rational if it is birational over \(k\) to \(\mathbb P_k^d\). Here \(k\)-unirational means that there exists a dominant rational map over \(k\) from \(\mathbb P_k^d\) to the variety.

\textbf{Questions, conditional on the intended ground field and curve--automorphism data.} Determine whether
\[
 X_{4,6}\dashrightarrow\mathbb P_k^4
 \quad\text{is a birational equivalence for some such map,}
\]
and determine whether there exists a dominant rational map
\[
 \mathbb P_k^5\dashrightarrow X_{5,6}.
\]
Equivalently, these ask whether \(Q_{4,6}\) is \(k\)-rational and whether \(Q_{5,6}\) is \(k\)-unirational. The missing field and automorphism specifications must be resolved before these conditional questions constitute a fully specified formulation of the source problem.
\subsection{Short English statement}
For the diagonal cyclic quotients of powers of the source's equianharmonic elliptic curve, determine whether the four-dimensional resolution is rational and whether the five-dimensional resolution is unirational. The ground field, distinguished automorphism, and higher-dimensional meaning of minimal resolution require clarification. These are statement-only questions requiring human mathematical review.
\subsection{Sources}
\begin{itemize}
\item Local UnsolvedMath registry, permanent ID \(30002829\), source problem code \texttt{catalogue entry}, title \emph{Rationality of Ueno-Type Varieties}; supplied complete source record in the Oberwolfach Reports Open Problems collection. The proposed year is recorded as 2015; no individual proposer is supplied.
\item \emph{Algebraic Geometry} (2015), Oberwolfach Reports, pp.\ 783--836, DOI 10.4171/owr/2015/15: \url{https://doi.org/10.4171/owr/2015/15}. This is also the supplied literature reference labelled OWR Report 15/2015.
\item Supplied literature reference: Catanese--Oguiso--Verra, \url{https://arxiv.org/abs/1506.01925}.
\item Supplied literature reference: Mellit, \emph{Ueno-type manifolds}, \url{https://doi.org/10.1080/10586458.2019.1691088}.
\item Literature assessment dated exactly \texttt{2026-08-21}: the supplied curation labels the problem partially solved, reports substantial unirationality and invariant-theoretic results, and states that rationality of \(X_{4,6}\) and unirationality of \(X_{5,6}\) remain unanswered. This is supplied dated curation, not an independently verified literature assessment; the record requests independent verification of time-sensitive status and cited claims.
\item The registry labels the statement \texttt{corrected\_verified} and reports that it was checked against the official Oberwolfach report. That verification claim is supplied provenance and has not been independently checked for this draft.
\end{itemize}
\end{document}
